\documentclass[11pt]{article}
\usepackage{handout}
\codigo{P3}
\titulo{Resultados em Fortaleza e Roteiro de Defesa}

\begin{document}
\fazertitulo

Este handout reúne o que o pipeline encontrou na malha real do corredor Benfica--Centro--Aldeota. Ele separa o que os resultados \emph{mostram} do que \emph{não mostram} e prepara a conversa com o Prof.\ José Soares. Os números vêm de \path{resultados/saida_fortaleza.txt}, gerado com:
\begin{lstlisting}[language=bash]
bash codigo/recortar_osm.sh ceara-latest.osm.pbf dados/osm/corredor.osm
python codigo/otimizacao_antes_depois.py --osm-xml dados/osm/corredor.osm
\end{lstlisting}
O mapa vem do OpenStreetMap (extrato do Ceará). As velocidades são \textbf{simuladas} pelo modelo do handout P1. Leia P1 e P2 antes. \pontos{12}

\section{A malha}

\begin{center}
\begin{tabular}{lc}
\toprule
Recorte (oeste, sul, leste, norte) & $(-38{,}548;\ -3{,}752;\ -38{,}490;\ -3{,}718)$\\
Cruzamentos $N$ / vias $E$ & 2470 / 4166\\
Grau médio, $\avg{k^2}/\avg k$ & 3,37; 3,66\\
Vias residenciais / terciárias & 2435 / 708\\
Vias secundárias / primárias / troncais & 207 / 789 / 18\\
Pontes no mapa completo & 177 (todas entradas de becos)\\
3-core & vazio\\
$p_c$ (percolação de ligações, $r$ uniforme) & $0{,}564\pm0{,}002$\\
\bottomrule
\end{tabular}
\end{center}

O mapa completo tem redundância: suas únicas pontes são entradas de becos (T3, Problema 1). Toda fragmentação relevante, portanto, vem do congestionamento.

\begin{center}
\includegraphics[width=0.95\textwidth]{fortaleza_fig4_mapa.png}
\end{center}
\textit{Um dia simulado. Cinza escuro: vias lentas ($r$ baixo). Azul: as 41 vias de maior betweenness, quase todas no corredor Padre Valdevino / Duque de Caxias / Domingos Olímpio. Laranja: as 41 ligações críticas escolhidas pelo algoritmo guloso nesse dia, espalhadas pela malha e em ruas locais.}

\section{A transição}

\begin{center}
\begin{minipage}{0.48\textwidth}\centering
\includegraphics[width=\textwidth]{fortaleza_fig1_transicao.png}
\end{minipage}\hfill
\begin{minipage}{0.48\textwidth}\centering
\includegraphics[width=\textwidth]{fortaleza_fig2_antes_depois.png}
\end{minipage}
\end{center}

\begin{itemize}
  \item $\avg{q_c} = 0{,}475\pm0{,}014$ (erro padrão), com desvio de $0{,}086$ entre dias. O pico de $\avg{S_2}$ chega a $\approx15\%$ de $N$.
  \item $\avg R = 0{,}489\pm0{,}004$, com desvio de $0{,}025$ entre dias. A robustez flutua muito menos que $q_c$ (relativamente, 5\% contra 18\%).
  \item Embaralhando as velocidades entre as vias, o esperado seria $F^{-1}(1-p_c) = 0{,}502$. As correlações custam $\approx0{,}03$ de resiliência (um efeito de $\approx2\sigma$; P1, Problema 7).
\end{itemize}

\section{Antes vs.\ depois}

\begin{center}
\includegraphics[width=0.95\textwidth]{fortaleza_fig3_orcamento.png}
\end{center}

Com $k=41$ vias ($0{,}98\%$ de $E$):
\begin{center}
\begin{tabular}{lcccc}
\toprule
& \multicolumn{2}{c}{tempo real} & \multicolumn{2}{c}{plano fixo}\\
Estratégia & $\Delta q_c$ & $\Delta R$ & $\Delta q_c$ & $\Delta R$\\
\midrule
Aleatória & $+0{,}002(1)$ & $+0{,}003(1)$ & $0{,}000(0)$ & $+0{,}004(1)$\\
Mais lentas & $+0{,}078(16)$ & $+0{,}026(3)$ & $+0{,}095(21)$ & $+0{,}030(4)$\\
Maior betweenness & $+0{,}079(21)$ & $+0{,}053(4)$ & $+0{,}089(25)$ & $\bm{+0{,}057(4)}$\\
Red bonds (uma vez) & $+0{,}141(14)$ & $+0{,}031(3)$ & $+0{,}026(11)$ & $+0{,}011(2)$\\
Ligações críticas & $\bm{+0{,}323(23)}$ & $\bm{+0{,}093(4)}$ & $+0{,}052(17)$ & $+0{,}023(4)$\\
\bottomrule
\end{tabular}
\end{center}
(O número entre parênteses é o erro padrão nos últimos dígitos: $0{,}323(23)$ significa $0{,}323\pm0{,}023$.)

\begin{ideia}
\textbf{As três conclusões, em ordem de confiança.}
\begin{enumerate}
  \item \emph{Em tempo real}, destravar as ligações críticas é de longe a melhor estratégia. Com menos de 1\% das vias, $q_c$ vai de $0{,}475$ para $0{,}798$. Isso é 4 vezes o ganho do betweenness e da força bruta em $q_c$, e 1,8 e 3,6 vezes, respectivamente, em $R$. (Diferença para o betweenness em $R$: $7{,}6\sigma$.)
  \item Os gargalos de percolação \emph{não} são as vias importantes do mapa: sobreposição zero com as 41 vias de maior betweenness.
  \item Os gargalos \emph{evoluem}: a sobreposição entre dias consecutivos é 0,033, cerca de 7 vezes o acaso, mas longe de 1. Num plano fixo, a vantagem desaparece, e quem vence é o betweenness ($\Delta R = 0{,}057$ contra $0{,}023$, $5{,}6\sigma$), porque o corredor de alta betweenness está travado \emph{todos} os dias.
\end{enumerate}
\end{ideia}

\section{Onde ficam as vias}

\textbf{Plano fixo, ligações críticas} (as 10 primeiras pela pontuação de Borda): Rua Teresa Cristina, Rua Joaquim Torres, Rua 25 de Março (dois trechos), Rua Antônio Pompeu, Rua Tomás Acioli, Rua Barão de Aracati, Rua João Brígido, Rua Henrique Rabelo e Rua General Clarindo de Queiroz. Quase todas são vias locais ou coletoras, com $\avg r$ entre 0,49 e 0,79.

\textbf{Maior betweenness} (41 trechos): Rua Padre Valdevino (26), Av.\ Duque de Caxias (8), Av.\ Domingos Olímpio (5) e Av.\ Aguanambi (2). Todos com $\avg r\approx0$ no modelo.

\begin{obs}
Não tire conclusões sobre ruas específicas sem dados reais de velocidade. O modelo concentra a carga num único corredor porque usa roteamento ``tudo ou nada'' (T3, Problema 6). O que os nomes mostram é o \emph{tipo} de via: os gargalos de percolação são ruas locais que ligam bolsões, não avenidas.
\end{obs}

\section{Sintético vs.\ real}

\begin{center}
\begin{tabular}{lcc}
\toprule
& sintética ($45\times45$) & Fortaleza\\
\midrule
$N$ / $E$ / $\avg k$ & 2025 / 3728 / 3,68 & 2470 / 4166 / 3,37\\
$\avg{q_c}$ / $\avg R$ & 0,622 / 0,603 & 0,475 / 0,489\\
$\Delta R$, críticas, tempo real & $+0{,}065$ & $+0{,}093$\\
$\Delta R$, betweenness, tempo real & $+0{,}017$ & $+0{,}053$\\
$\Delta R$, críticas, plano fixo & $+0{,}023$ & $+0{,}023$\\
$\Delta R$, betweenness, plano fixo & $+0{,}018$ & $+0{,}057$\\
Jaccard entre dias consecutivos & 0,019 & 0,033\\
\bottomrule
\end{tabular}
\end{center}

\section{O que o trabalho \emph{não} mostra (diga antes que perguntem)}

\begin{enumerate}
  \item \textbf{As velocidades são simuladas.} A topologia é real, mas $r_e$ vem de um modelo calibrado à mão. As conclusões são sobre o \emph{mecanismo}, não sobre a Fortaleza de hoje. Próximo passo: GPS dos ônibus da Etufor ou dados da AMC.
  \item \textbf{Grafo não-direcionado.} Usamos conectividade fraca. Com mão única, o certo é usar componentes fortemente conexos.
  \item \textbf{Sem reroteamento.} A carga vem da betweenness em fluxo livre, sem equilíbrio de Wardrop. Isso exagera a concentração nos corredores.
  \item \textbf{Destravar = $r\to1$} é um cenário idealizado (o melhor caso para todas as estratégias).
  \item \textbf{Um único recorte}, sem análise de tamanho finito.
  \item \textbf{O guloso não é ótimo.} Não sabemos quão longe ele está do ótimo.
\end{enumerate}

\section{Roteiro de defesa: perguntas prováveis}

\textbf{``Por que definir $q_c$ pelo pico de $S_2$?''} Porque $S_2$ faz o papel da susceptibilidade: tem pico onde os dois maiores clusters, ambos fractais, estão prestes a se fundir (T1, Seção 4). Além disso, o pico é calculado \emph{exatamente} numa única varredura, junto com a ligação que o causa (T1, Exemplo 3). Um critério como $S_1 = N/2$ depende de uma escolha arbitrária e não identifica nenhuma ligação.

\textbf{``O pico de $S_2$ não deveria sumir no limite termodinâmico?''} A fração $S_2^{\max}/N\sim L^{d_f-2} = L^{-5/48}$ vai a zero, mas devagar demais para importar em qualquer cidade real: para $L=50$, ainda é $\approx0{,}66$ vezes a amplitude.

\textbf{``Você achou centenas de pontes, mas Coniglio diz $L^{3/4}$.''} A maioria das pontes segura galhos (\emph{dangling ends}). As red bonds de Coniglio são as pontes do \emph{backbone}, as de massa de corte macroscópica. Essas crescem como $L^{0{,}62}$ no teste, perto de $L^{3/4}$ (T1, Seção 5).

\textbf{``Seu $q_c$ às vezes cai depois de uma melhoria. O código está errado?''} Não. $q_c$ é a posição de um máximo e pode pular entre eventos de fusão (T1, Problema 7). Por isso reportamos também $R$, que é provadamente monótono (testado em 200 tentativas).

\textbf{``Por que as ligações críticas não são as de maior betweenness?''} Porque a criticidade é uma propriedade do grafo \emph{congestionado}. No mapa completo, a ligação crítica é uma rua comum, com desvios. Quando os desvios estão lentos, ela vira ponte entre dois clusters macroscópicos, com betweenness $\sim S_1S_2$ no grafo congestionado (T3, Ideia 4).

\textbf{``Se o plano fixo não funciona, para que serve o método?''} Para o controle adaptativo. O cálculo leva 4 ms. A CTAFOR poderia recalcular as ligações críticas a cada poucos minutos com dados de GPS e priorizar o verde nelas. Para obras e planos fixos, a resposta do modelo é atacar os corredores cronicamente saturados.

\textbf{``Qual é a classe de universalidade?''} A malha é plana e as correlações do modelo são de curto alcance ($\sim2$--3 quarteirões), então esperamos percolação 2D ($\nu=4/3$, $d_f=91/48$). Ainda não verificamos. O teste seria uma análise de tamanho finito com recortes de tamanhos diferentes (Problema 2).

\problema{2} Em vias arteriais de 60 km/h, $q_c = 0{,}475$ corresponde a $28{,}5$ km/h. Dê uma interpretação concreta para um motorista: o que significa ``a cidade se fragmenta em $q_c$''? O que muda, nessa linguagem, quando $q_c$ sobe para $0{,}798$?

\problema{4} \textbf{Desenhe o próximo experimento: tamanho finito em Fortaleza.} Usando recortes concêntricos de lados $L$, $2L$ e $4L$ (por exemplo, centrados na Praça do Ferreira):
\begin{enumerate}
  \item Que grandezas você mediria em cada recorte, e que leis de potência espera para cada uma?
  \item Como estimaria $\nu$ e $d_f$ a partir delas?
  \item Cite dois motivos pelos quais uma cidade real pode \emph{não} obedecer a essas leis (pense na heterogeneidade dos bairros e no tamanho máximo possível).
\end{enumerate}

\problema{3} Por que o betweenness funciona muito melhor em Fortaleza ($\Delta R = 0{,}053$ em tempo real) do que na malha sintética ($0{,}017$)? Compare a estrutura das avenidas nas duas malhas.

\problema{3} \textbf{Falsificação.} Proponha um teste que poderia \emph{derrubar} a conclusão 1 da Ideia 1. Por exemplo: o que aconteceria se a intervenção fosse mais realista, $r_e\to r_e+0{,}2$ em vez de $r_e\to1$? E se o orçamento fosse 0,2\% em vez de 1\%? Diga qual resultado você esperaria e qual resultado o deixaria preocupado.

\solucoes

\sol{1}
$q$ é o nível de exigência. Em $q=0{,}475$, se o motorista só aceita vias que entregam pelo menos 47,5\% da velocidade livre (28,5 km/h numa arterial, 14 km/h numa rua de 30), a cidade se parte em ilhas: não existe rota ``aceitável'' entre a maioria dos pares de bairros. Com $q_c = 0{,}798$, o mesmo motorista conseguiria exigir 80\% da velocidade livre (48 km/h numa arterial) e ainda atravessar a cidade. Na prática: existem rotas fluidas entre quase todos os pares de pontos, mesmo no pico.

\sol{2}
(a) Para cada recorte: $\avg{q_c}(L)$ e a largura $\delta q_c(L)$ (desvio entre dias); $\avg{S_2^{\max}}(L)$; o número de pontes de massa macroscópica no ponto crítico, $L_{\rm red}(L)$. As previsões são $\avg{q_c}(L)-q_c(\infty)\sim L^{-1/\nu}$, $\delta q_c\sim L^{-1/\nu}$, $S_2^{\max}\sim L^{d_f}$ e $L_{\rm red}\sim L^{1/\nu}$. (b) Inclinações em gráficos log-log: $\delta q_c$ e $L_{\rm red}$ dão $1/\nu$ (duas estimativas independentes, uma para conferir a outra), e $S_2^{\max}$ dá $d_f$. Confira a hiperescala $d_f = 2-\beta/\nu$. (c) Os bairros têm morfologias diferentes (grade regular na Aldeota, traçado irregular no Centro), então aumentar $L$ muda o \emph{tipo} de malha, não só o tamanho. E a cidade é finita: só cabem 2 ou 3 décadas de $L$, pouco para leis de potência limpas. Correções de escala serão grandes.

\sol{3}
A malha sintética tem avenidas regulares a cada 9 quadras, todas equivalentes: a carga se divide entre várias avenidas paralelas, e destravar 1\% das vias de maior betweenness só abre pedaços de várias delas. Em Fortaleza, poucos corredores longos atravessam o recorte (Padre Valdevino, Duque de Caxias, Domingos Olímpio), e o roteamento ``tudo ou nada'' concentra a carga neles. Destravar o corredor inteiro cria uma ``espinha'' rápida que conecta muitos bolsões em qualquer $q$. Parte desse efeito é real (Fortaleza tem mesmo corredores dominantes), e parte é artefato da falta de reroteamento.

\sol{4}
Com $r_e\to r_e+0{,}2$, todos os ganhos caem. A pergunta é se a \emph{ordem} das estratégias se mantém. Esperamos que sim: as ligações críticas estão perto de $q_c$ ($\avg r\approx0{,}5$--0,8), então subir 0,2 já as tira da zona crítica, enquanto as vias mais lentas ($r\approx0$) continuariam fechadas com $+0{,}2$. Com 0,2\% de orçamento, esperamos uma vantagem relativa ainda \emph{maior} das ligações críticas, já que o ganho delas é concentrado nas primeiras vias (Fig.\ 3). Ficaria preocupante se, com intervenção parcial, as ligações críticas perdessem para o betweenness também em tempo real: isso indicaria que o ganho vem de abrir ``atalhos perfeitos'' irreais, e não de entender o ponto crítico.

\end{document}
